\section{Pipeline Overview}

The project is a single pipeline organised into three stages, from a generic pretrained
model to the system presented here. Each stage adds a capability to the previous
checkpoint without discarding what was already acquired (Figure~\ref{fig:arc}).

\begin{figure}[h]
\begin{verbatim}
STAGE A1   domain adaptation
   SmolLM2-360M continued-pretrained on ~12k Djinni IT postings (LoRA)
   in-domain PPL 16.37 to 11.38 (-31%); out-of-domain unchanged
   yields a domain-adapted 360M checkpoint (merged LoRA)

STAGE A2   alignment
   SFT on ~7.5k request-posting pairs (validation PPL 11.66 to 4.54)
   DPO via RLAIF, beta = 0.10; win-rate 8-0 vs SFT; DPO val PPL 6.01
   yields ma2-360m-dpo-b01, the Job-Offer specialist (~430 tok/s)

STAGE A3   the system (this work)
   gemma-4 E4B orchestrator + ma2-360m-dpo-b01 (live Job Offer drafter)
   Component 1  Hybrid RAG: vector + graph, 10 domains + 210,048 CVs
   Component 2  Agentic pipeline: 8 LLM-decided steps, memory, guard
   Component 3  LLM-as-judge: RAGAS faithfulness + answer-relevance
   Component 4  Performance optimisation: llama.cpp, router, slot cap
   Supporting   Chain-of-thought reasoning, in the Thinking panel
\end{verbatim}
\caption{The three-stage arc from raw pretrained model to job2cool. The A3 stage is not a
black box: it expands into the live, health-monitored dependency graph the system renders
of itself (Section~\ref{sec:infra}).}
\label{fig:arc}
\end{figure}

\subsection{A1: domain adaptation}

Mini-Assignment~1 continued the pretraining of SmolLM2-360M on roughly 12,000 IT job
postings from the Djinni dataset (141,897 postings in the full release), shifting the
model's language distribution toward the vocabulary and conventions of this type of
posting. Under identical conditions, LoRA beat full fine-tuning on in-domain perplexity
(11.38 versus 13.12, against a base of 16.37) while training only 2.3\% of parameters in
roughly a third of the time. Out-of-domain perplexity, measured on LinkedIn postings,
remained practically unchanged for both, confirming the adaptation was IT-specific rather
than a general drift. The merged LoRA checkpoint became the A2 starting point.

\subsection{A2: alignment}

Mini-Assignment~2 turned that domain-fluent but instruction-blind checkpoint into a
recruiter-instructable drafter. Supervised fine-tuning on roughly 7,500
request-and-posting pairs drove validation perplexity from 11.66 to 4.54 and made the
model reliably emit the four Job-Offer sections (Summary, Required Skills,
Responsibilities, and Requirements) in Markdown. DPO via Reinforcement Learning from AI
Feedback then ranked sampled candidates against a rubric covering faithfulness, structural
completeness, language quality, and absence of repetition; a sweep over four values of
$\beta$ peaked at $\beta = 0.10$, yielding \verb|ma2-360m-dpo-b01|, which won 8 of 20
held-out comparisons against the SFT model (0 losses) under a Granite judge with
order-swap. DPO validation perplexity settled at 6.01, higher than SFT's 4.54, as
expected under KL regularisation.

Two A2 findings carry directly into A3. First, the aligned model's behaviour depends
strongly on the prompt template and on inference parameters, most notably the repetition
penalty, whose absence caused repetition collapses, so A3 must keep the same inference
rigour when calling the model through a new orchestration layer. Second, judge-based
evaluation has a discrimination ceiling: a smaller judge distinguished broad comparisons
(base versus aligned) well but was unreliable on close comparisons between similar
checkpoints, which motivates an evaluation design that does not rest on a single judge
(Section~\ref{sec:eval}).

\subsection{A3: the system, and the model-choice justification}

Mini-Assignment~3 uses \verb|ma2-360m-dpo-b01| exactly as delivered, without further
modification, assigning it one specific role: drafting Job Offers. Everything else, namely
interpreting the conversation, resolving the knowledge-base domain, deciding which
documents to produce, retrieving evidence, composing each section (including refining the
A2 draft), and attaching citations, is handled by the orchestration layer and the
\verb|gemma-4| general model.

The assignment permits switching to a different open-source model or running a further
adaptation pass. job2cool does neither as a replacement: it retains the A2 model and adds
\verb|gemma-4| alongside it, and the justification is capability-versus-task fit. The A2
model is a 360M-parameter specialist, excellent and fast at the one task it was aligned
for (natural-language need into a Markdown Job Offer), but with no tool use, no
long-context planning, and no general composition ability for onboarding, interview, or
culture content; asking a 360M model to drive the whole pipeline would fail. \verb|gemma-4|
is the open-source orchestrator (4B-class, 128\,K context) that runs the agentic plan,
composes the RAG-grounded sections, and refines the A2 draft. Both are co-resident on
\verb|agent_server| and selectable per request by model id, so the specialist and the
orchestrator coexist rather than one superseding the other. This is the honest reading of
the arc: the final stage is an orchestration-and-grounding layer that makes the aligned
model \emph{usable}, not a larger model that makes it redundant. Whether \verb|gemma-4|'s
composition step masks the A2 alignment in the final output is one of the questions the
evaluation is built to answer (Section~\ref{sec:eval}, Section~\ref{sec:arc}).

Retrieval and orchestration therefore operate on a different axis from the earlier work,
which improved the quality of one output (the Job Offer) in isolation. They address what
an isolated aligned model cannot resolve regardless of its quality: knowledge of the
company's actual material, determining what a request is asking for, holding context
across a conversation, and proving through citation that a claim has an identifiable
source.
